\section{Morphosyntactic transparency}\label{sec:2}

Morphosyntactic transparency is the case where $X$ is a syntactic construction (or a lexical class within one) and $R$ is a grammatical relation sensitive to a property borne inside $X$ but not by its head. Feature matching between a non-head property and a co-occurring exponent is the paradigm, and agreement is the paradigm of that. Selection, category licensing, binding, and dependency formation are possible targets, as §\ref{sec:2-5} shows for transparent free relatives.\footnote{Extraction transparency (whether an island or barrier is opaque or transparent to long-distance dependency formation) is a morphosyntactic-accessibility relation under the schema (with $R$ = dependency formation), and the associated question remains active. Transformational accounts posit empty elements (traces, copies) that mediate the connection between a fronted element and its base position; non-transformational dependency accounts \autocite{gibson2025,dacunha-gibson-2025-syntactic-complexity} argue that no such mediating empty elements are needed and that direct head-to-dependent connections better fit acceptability and processing data. Both sides make schema-level claims about whether $X$ contains invisible mediating structure. The centre of the argument here is the agreement-style cases.} Subject-verb agreement is the textbook instance. The concrete contrasts are simple. In (\ref{ex:cgel-two-direction}a), agreement tracks \mention{spots} through singular \mention{number}; in (\ref{ex:tfr-r-relations}a), agreement tracks \mention{trousers} through singular \mention{what}.

This section carries the paper's primary field-relative claim. The projected targets are syntactic and morphosyntactic: agreement, syntactic-category licensing, complement selection, behaviour under complement ellipsis, and distributional restrictions inside NP and fused-relative constructions. Nothing in the syntactic result depends on phrase-structure representation as such: the profiles require only observable relations among a mediator, an internal property, and the grammatical relation that tracks it, whether those relations are represented in constituency or dependency terms. The section contains two grammatical cases, but only the QN section carries a prospective candidate-member test; the TFR result remains a conditional illustration based on Kim's constructional analysis and pilot corpus evidence.

\subsection{Number-transparent quantificational nouns}\label{sec:2-1}

The paradigm again is \textit{CGEL}'s two-direction contrast in (\ref{ex:cgel-two-direction}): plural agreement despite singular \mention{number}, and singular agreement despite plural \mention{heaps}. Schema-wise (under the umbrella schema), $X$ is the number-transparent quantificational noun + \mention{of}-complement construction; $P$ is the number/count profile of the \mention{of}-complement; $R$ is subject-verb agreement.

\textit{CGEL}'s \enquote{simple agreement rule} is head-driven: \enquote{the verb agrees with a subject with the form of an NP whose person--number classification derives from its head noun} (§18.1, p.~499), as in (\ref{ex:simple-agreement}).

\ea\label{ex:simple-agreement}
\ea \mention{The nurse wants to see him.}
\ex \mention{The nurses want to see him.}
\z
\z

Number-transparent NPs override the rule. \textit{CGEL} describes the override as obligatory in the central cases; rare singular-verb attestations with \mention{number} are noted in footnote 73 to §18.2 (p.~502) and analyzed there as hypercorrections rather than as variant grammar. The contrast with the optional override of collective nouns (§\ref{sec:2-2}) is the operative point. The construction's job is to let the \mention{of}-complement's number/count profile reach agreement past the head's morphological number, which would otherwise have determined it.

The core inventory is small and well documented. \textit{CGEL} identifies the principal members in example [57] (§3.3, p.~350): \mention{lot, plenty} (singular forms; \mention{lot} most often with \mention{a}-determiner, \mention{plenty} typically without a determiner); \mention{lots, bags, heaps, loads, oodles, stacks} (plural forms, typically without a determiner); \mention{remainder, rest} (singular forms typically with \mention{the}-determiner); \mention{number, couple}₂ (singular forms typically with \mention{a}-determiner, taking only plural \mention{of}-complements, partitive or pseudo-partitive). Of these, \enquote{the clear cases of number-transparent singular nouns are \mention{lot, number}, and \mention{couple}₂} (p.~503); \mention{majority} and \mention{minority} are borderline: plural override is \enquote{surely obligatory} in examples like \mention{The majority of her friends are Irish}, but singular agreement is occasionally attested (§18.2, pp.~503--504). The determiner patterns above are typical, not absolute: \mention{the lot of it} (in a totality sense) is grammatical alongside \mention{a lot of it}, and similar sense-conditioned variations exist for other members.

Keizer's \autocite*{keizer-2007-english-noun-phrase} study of English NP categorization and Van Eynde and Kim's \autocite*{van-eynde-kim-2023-pseudo-partitives} Head-driven Phrase Structure Grammar (HPSG) analysis of pseudo-partitives are the immediate background for treating these strings as structured binominal configurations rather than as a single loose family. Within the binominal family, Kim \autocite*[98--103]{kim-2026-form-function} distinguishes three sub-types by agreement head, illustrated in (\ref{ex:three-subtypes}): measure nouns (\mention{pound}), where the first noun controls agreement; quantity nouns (\mention{lot}, \mention{number}), where the second noun controls it; and collective nouns (\mention{group}, \mention{bunch}), which allow either.

\ea\label{ex:three-subtypes}
\ea Measure (N1 controls): \mention{One pound of tomatoes equals 90 calories.}
\ex Quantity (N2 controls): \mention{A lot of marbles are about to fall on the floor.}
\ex Collective (either): \mention{A group of police cruisers is speeding across the bridge}; \mention{a group of immigrants move in}.
\z
\z

The three example sentences are Kim's, from COCA \autocite[99--102]{kim-2026-form-function}. The quantity type is \textit{CGEL}'s number-transparent class, and the measure type gives the schema its concealing variant: the same binominal string with agreement held by the first noun (§\ref{sec:2-3}).

Class membership predicts a cluster of behaviours beyond the agreement override. The relevant contrast is the standard \term{partitive} / \term{pseudo-partitive} distinction: partitive \mention{of}-complements require a definite complement that picks out a subset of a definite set (\mention{of the money}, \mention{of those books}, \mention{of these soldiers}), while pseudo-partitive \mention{of}-complements express a quantity, measure, or collection relation with a bare or otherwise non-definite complement in the QN frame (\mention{of money}, \mention{of soldiers}, \mention{of some paintings}) \autocite{selkirk-1977-noun-phrase-structure,stickney-2007-pseudopartitive-partitive,kim-michaelis-2020-syntactic-constructions,van-eynde-kim-2023-pseudo-partitives}.\footnote{The headedness issue is framework-internal. HPSG analyses treat the first noun as syntactic head in \mention{A of the B} partitives, while \textit{CGEL} retains the quantificational noun as head and treats agreement with the complement as override. Since the schema is framework-portable rather than neutral (§\ref{sec:3}), the paper doesn't adjudicate the headedness choice; it records that the verdict is conditional on it.}

\mention{Lot} most often takes \mention{a} as determiner with limited premodification (\mention{A whole/huge lot of time has been wasted}; \textit{CGEL}, §3.3, p.~350); \mention{the lot of it} is also available in a totality sense, but other determiners are blocked (*\mention{some lot of it}). \mention{Plenty} typically takes no determiner or modifier (\mention{plenty of butter}, not generally \mention{a remarkable plenty of butter}), and like the plural-only forms it strongly prefers the bare \mention{of}-complement.\footnote{A COCA pilot (May 2026, extended July 2026) supports the preference but not a rate. The bare complement dominates for \mention{plenty} and the plural-only forms, while \mention{rest} and \mention{remainder} prefer the determined one on matched complements ($\approx\!97\%$ for \mention{day, time, people, money}). No overall share is defensible: the two frames were sampled over different complement inventories, and frequency in any case supplies no denominator of occasions on which a partitive reading was intended. The partitive frame is attested for every member of both subgroups. Full counts and methodology in Appendix~\ref{a.2-partitive-vs.-pseudo-partitive-of-complements}; raw data in the supplementary material (see Data availability).} The plural-only forms (\mention{lots, bags, heaps, loads, oodles, stacks}, and \mention{piles} beyond \textit{CGEL}'s list) show the same pattern.

\ea\label{ex:plural-only-of-complements}
\ea \mention{They have oodles of money.}
\ex \mention{Oodles of the money had already been spent.}
\z
\z

The contrast in (\ref{ex:plural-only-of-complements}) shows the plural-only pattern: the first is the characteristic use, the second a marked one. It is not a blocked one. A wildcard sweep of the determined frame attests it for every member of the subgroup, at 465 tokens across 374 distinct complements, where a probe over four hand-picked complements had returned 13 and two zeros (Appendix~\ref{a.2-partitive-vs.-pseudo-partitive-of-complements}). \mention{Remainder} and \mention{rest} are a different case again; \mention{number} and \mention{couple}₂ take plural obliques.

\ea\label{ex:partitive-plural-oblique}
\ea \mention{the remainder of the time}
\ex \ungram{\mention{the remainder of time}} (partitive reading only; fine as \enquote{all time henceforth})
\ex \mention{the rest of society}
\ex \mention{a number of the protesters}
\ex \ungram{\mention{a number of money}}
\z
\z

The obliques in (\ref{ex:partitive-plural-oblique}) are structured, but the lexical classification of \mention{rest} is not settled by the corpus pattern. Under the semantic criterion adopted here, \mention{rest} and \mention{remainder} denote the complement subset of a presupposed whole rather than a quantity of it, so the pseudo-partitive construction, defined as one in which N1 denotes \enquote{a quantity or amount} of what N2 denotes \autocite[271]{van-eynde-kim-2023-pseudo-partitives}, is unavailable. A competing analysis could treat bare \mention{the rest of society} as pseudo-partitive by treating \mention{rest} as a part or amount expression. The data establish the bounded-and-unique condition on the bare frame, but do not decide that lexical-semantic question; Table~\ref{tab:qn-matrix} therefore marks \mention{rest} as unresolved while retaining the semantic analysis as the working description.

For the bare-complement reading, the whole must be uniquely identifiable \mention{and} bounded, so that it has a proper part to take a remainder of. Generic collectives supply both unaided, which is why (\ref{ex:partitive-plural-oblique}c) needs no determiner and is common: \mention{the rest of society} alone has 442 COCA tokens, with \mention{humanity} 248, \mention{nature} 78, \mention{mankind} 77, \mention{creation} 42, and \mention{Congress} 21 behind it. An unbounded complement supplies neither. Bare \mention{of time} offers no partition, which is why (\ref{ex:partitive-plural-oblique}b) is starred on the relevant reading and why the string survives only on the unrelated duration reading, \enquote{all time henceforth}. That reading accounts for most of the corpus tokens of \mention{the rest of time}, so those tokens are not evidence about the partitive/pseudo-partitive classification (Appendix~\ref{a.2-partitive-vs.-pseudo-partitive-of-complements}).

\textit{CGEL}'s own minimal pair for \mention{lot} has the complement controlling agreement in both directions (p.~349).

\ea\label{ex:override-visibility}
\ea \mention{A lot of work was done.}
\ex \mention{A lot of errors were made.}
\z
\z

The pair in (\ref{ex:override-visibility}) shows what number transparency \mention{is} in the umbrella schema's vocabulary. $X$ (the construction) leaves $P$ (the number/count profile of the \mention{of}-complement) accessible to $R$ (agreement). What $X$ conceals from agreement is the morphological number of the head noun \mention{lot, plenty, heaps}, etc. The agreement-targeting feature is on the \mention{of}-complement, not on the surface head. That's what makes $X$ transparent with respect to $P$: $X$ lets $P$ reach the verb anyway.

The override is only visible where the head's number and the complement's diverge. (\ref{ex:override-visibility}b) shows it (singular \mention{lot}, plural \mention{errors}); (\ref{ex:override-visibility}a) doesn't, since head and complement both call for the singular. The discriminating cases are the crossed ones: a singular-form quantificational noun with a plural complement, or a plural-form one (\mention{lots, heaps}) with a non-count complement. The matched cases are silent, not counterevidence. Transparency is therefore under-detected: a class can be transparent while most of its tokens look head-driven, so corpus frequency understates it.

\subsection{Collective NPs as a transitional case}\label{sec:2-2}

Singular collective nouns (\mention{committee, jury, family, team, government}; also \mention{clergy, intelligentsia, public}, etc.) allow optional plural override of head-driven agreement, as in \textit{CGEL} [8] and [10] (§18.2, pp.~501--502).

\ea\label{ex:collective-agreement}
\ea \mention{The committee has}/\mention{have} \mention{not yet come to a decision.}
\ex \mention{The committee consists}/\ungram{\mention{consist}} \mention{of two academic staff and three students.}
\ex \mention{One committee, appointed last year, has}/\marg{\mention{have}} \mention{not yet met.}
\z
\z

The override in (\ref{ex:collective-agreement}) is construal-sensitive: singular agreement tracks the unit construal, plural the member construal. Plural is unavailable with predicates that apply to the whole but not individual members, and is \enquote{of questionable acceptability} if the collective takes determiner \mention{one} or \mention{another} (p.~502).

Both collective NPs and number-transparent quantificational NPs override the head's morphological singular at agreement, and the override is possible because plural agreement requires only non-singular construal, not the exact enumeration that singular \mention{a}(\mention{n}) and low cardinals demand. That is why tight determiners degrade the override that (\ref{ex:collective-agreement}a) allows while plural agreement survives it: the override operates at the loose end of the count cluster, not across it. Reynolds \autocite*{reynolds2026countability} develops that loose/tight dissociation at length.

The two cases differ in where the non-singular construal comes from. In collectives it is lexically encoded in the head, so the override is selective: it needs a predicate that licenses member-level construal, which is why whole-unit predicates block it. In number-transparent QNs it is imported from the \mention{of}-complement, so the override applies regardless of predicate type. That difference matters more than it first appears, because only one of them satisfies the schema as stated. The schema requires one element within $X$ to bear $P$ and another to conceal it. In the QN case those are distinct expressions: the \mention{of}-complement bears the number profile, the quantificational head conceals it. In \mention{the committee have decided} they are not. The singular morphology and the member-level plurality are two representations of one lexical head, so there is no second element doing the concealing. Collectives are therefore a near case: a competition between morphological and semantic properties within a single sign, not access across a mediator. Broadening \mention{element} to cover representational layers would accommodate them, at the cost of the non-instance test's grip, since almost any property could then be said to conceal another representation of itself. The transitional label in this subsection's title is meant in that strict sense.

Alongside \textit{CGEL}'s override framing, the HPSG tradition encodes the same observation as index rather than concord agreement: concord tracks morphosyntactic features internal to the NP, while index agreement can track the referential or semantic index available to the agreement target \autocite{wechsler-zlatic-2003-agreement,kim-2004-hybrid-agreement}. Kim and Sells \autocite*[64--65]{kim-sells-2015-binominal-nps} run the two levels on the cases at issue here: measure \mention{four pounds was} is morphosyntactically plural but index-singular, while collective \mention{this team have} is the mirror image, morphosyntactically singular with a plural index. On that encoding, collective agreement makes semantic plurality available to agreement, but it isn't access transparency in the strict sense used here, because the bearer and the proposed mediator aren't distinct elements. The index/concord split is consistent with the framework-portability of §\ref{sec:3}.

\textit{CGEL} itself notes that \enquote{the division between collective and number-transparent nouns is by no means sharply drawn} (§18.2, p.~502). Some nouns straddle. \mention{Couple} has two senses: \mention{couple}₁ (a man and woman in a relationship) is a regular collective; \mention{couple}₂ (`approximately two') is number-transparent. \mention{Bunch} heads the same binominal \mention{of}-frame as §\ref{sec:2-1}'s quantity nouns but is a collective noun in Kim's \autocite*[98--103]{kim-2026-form-function} classification, the sub-type where either noun can control agreement; the straddling \textit{CGEL} reports is what that classification predicts. \mention{Bunch} is also the case where the profile framework does demotion work: \textit{CGEL} claims it tilts collective for inanimate aggregates but transparent for groups of people, illustrated by the contrast in (\ref{ex:bunch-transition}).

\ea\label{ex:bunch-transition}
\ea \mention{A bunch of flowers was presented to the teacher.}
\ex \mention{A bunch of hooligans were seen leaving the premises.}
\z
\z

\textit{CGEL} treats the first as collective and the second as transparent (§18.2, p.~503). A pilot COCA check (May 2026) supports the animate half, where \mention{bunch}-subjects are consistently plural, and finds no positive support for the inanimate half: inanimate \mention{bunch} is rare in subject position, \textit{CGEL}'s own example is unattested, and the few tokens that surface take plural rather than the predicted singular. Corpus absence of a constructed example isn't ungrammaticality, so this doesn't touch \textit{CGEL}'s contrast as a claim about grammar; what it shows is that the inanimate-collective pattern doesn't project as a corpus regularity. That localizes the projectibility claim rather than the grammatical one, and marking where a textbook contrast doesn't project is what demotion conditions are for. Counts in Appendix~\ref{a.4-bunch-animate-vs.-inanimate-aggregates}.

The paradigm and the transitional case also differ in internal structure. Number-transparent NPs are class-like, with the override the default behaviour for class members. Collective NPs are gradient: the override probability depends on the head, the predicate, the determiner, the dialect register, and the construal the speaker has in mind. A theory of morphosyntactic transparency that ignored the difference between a class-based default override and a construal-driven optional one would lose the interesting structure of both cases.

\subsection{Coordination as a non-instance}\label{sec:2-3}

Coordination of singulars triggers plural agreement.

\ea\label{ex:ordinary-coordination}
\mention{Bob and Jane are happy.}
\z

The ordinary coordination in (\ref{ex:ordinary-coordination}) looks like a transparency case at first (the conjuncts' number features percolate up to the coordinate construction and determine agreement), but the schema excludes it. The schema requires that $X$ \mention{could have} made $P$ inaccessible. The coordinate construction does no such hiding: there's no head whose morphological number competes with the conjuncts' features for agreement.\footnote{This verdict presupposes a headless (or numberless-head) analysis: with no conjunct serving as a number-bearing head, nothing conceals the conjuncts' summed features. The picture changes under an analysis on which one conjunct is the head. Closest-conjunct agreement \autocite{bhatt-walkow-2013} is the case in point: where a verb tracks the nearest conjunct, that agreement conceals the resolved features and resolved agreement reaches them past it, so the same test can return an instance. As with the language-internal concealing baseline of §\ref{sec:6}, the verdict rides on the prior syntactic analysis.} The plural agreement is compositional addition (1 + 1 = 2), not mediation.

Some coordinate NPs take singular agreement.

\ea\label{ex:holistic-coordination}
\ea \mention{Bacon and eggs is my favourite breakfast.}
\ex \mention{Wine and cheese is expected.}
\ex \mention{Fish and chips is expensive these days.}
\z
\z

Some of the same coordinate strings also take plural agreement.

\ea\label{ex:coordination-plural-counterparts}
\ea \mention{Wine and cheese make the perfect party pair.}
\ex \mention{Fish and chips taste their best right after frying.}
\z
\z

The singular holistic cases in (\ref{ex:holistic-coordination}) and the plural counterparts in (\ref{ex:coordination-plural-counterparts}) leave the exclusion intact. In such cases agreement tracks a construal of the coordinate phrase: either a dish, package, event type, or conventional pairing, or a plural set of jointly mentioned items. The singular feature isn't a property of either conjunct that remains accessible through the coordinate construction; it's a construal of the coordinate phrase as a whole. Ordinary plural coordination and holistic singular coordination therefore differ in agreement semantics, but neither is transparent in the schema's sense. The former is compositional summation; the latter is whole-phrase unit construal.\footnote{Cases like \mention{One and one is two} are an arithmetic-expression construction, not a meal/package coordination: \mention{and} functions close to an addition operator, and the subject denotes an operation, equation, or expression. Not transparency either.} \textit{CGEL} treats coordination-and-agreement under a separate heading (§18.4), not alongside the override constructions of §18.2.

The clean distinction across §\ref{sec:2}'s cases: in number-transparent QNs, complement number becomes accessible through a quantificational head; in collective plural agreement, member plurality is available to agreement but is borne by the head itself, so no distinct mediator intervenes; in unit-construal coordination, the whole coordinate phrase is construed as singular and no internal property is being made transparent. Only the first is access through a mediator.

The \enquote{$X$ could have made $P$ inaccessible} clause works as an independent diagnostic, not an ad hoc filter. The test asks whether the construction-type admits a concealing counterpart: is there a version of the same construction-type, attested in the language, in which $X$'s internal structure makes $P$ inaccessible to $R$? For English NPs, head-driven agreement is the default and the number-transparent class is the marked exception within the same NP type. For coordinate constructions no such counterpart exists: the type admits no version in which conjunct features fail to reach agreement, so coordination has no control against which access can be measured. The diagnostic distinguishes mediator-types from compositional-summation types: access is informative only where the same mediator type has opaque or partially opaque alternatives.

One condition has to be attached, or \mention{same construction-type} becomes adjustable after the fact. Drawn narrowly enough, no transparent construction has an opaque sibling and every case is a non-instance; drawn broadly enough, almost anything does and the filter passes everything. So the concealing control must belong to a constructional family that is independently motivated, established on criteria that don't appeal to the transparency effect the test is assessing. English head-driven agreement qualifies: it is the unmarked pattern for NP subjects on grounds that have nothing to do with quantificational nouns, which is what makes the number-transparent class a marked exception \mention{within} a family fixed elsewhere. Tsez class IV agreement with a clausal argument qualifies on the same footing. Where no such family can be identified without invoking the effect, the test returns no verdict rather than a negative one; the typological cases in §\ref{sec:6} show why that control has to be fixed language-internally.

The condition also locates where the diagnostic is soft. Deciding what counts as the same construction-type is itself a language-internal analysis, so the test is crisp where the family is independently motivated and loose where it isn't. That is a real limitation on the paper's central instrument, not a detail about its edges.

A second non-instance fails on a different role, and the contrast is worth having because it answers the obvious deflationary reading of the QN case. In deferred or metonymic reference, agreement can track a referent the sentence never names.

\ea\label{ex:reference-transfer}
\ea \mention{The hash browns at table nine are}/\ungram{\mention{is}} \mention{getting cold.}
\ex \mention{The hash browns at table nine is}/\ungram{\mention{are}} \mention{getting angry.}
\z
\z

Kim's \autocite*[1108]{kim-2004-hybrid-agreement} pair, which he credits to Nunberg, turns on which entity the subject is used to pick out: the food in the first, the customer who ordered it in the second. Warren \autocite*[16]{warren2006} states the generalization, that metonymic subjects \enquote{need not agree as to number with their predicates} where metaphorical ones consistently do, and adds that a metonymic pronoun \enquote{sometimes agrees with the explicit and sometimes with the implicit element of the expression}, the choice settled by context.

Filled into the frame of §\ref{sec:1}, the \textit{bearer} cell is empty. The singular in (\ref{ex:reference-transfer}b) is sourced from the designated customer, who is expressed nowhere in the subject NP, so there is no element within $X$ bearing $P$ and nothing for the relation to reach \mention{through}. Coordination fails the test for want of a mediator; this fails for want of a bearer. Both are outside the schema, for different reasons.

Number transparency can be read as construal-tracking: the verb agrees with how the speaker conceives the subject, rather than with any constituent. Reference transfer shows what construal-tracking looks like on its own, and it looks different. It's context-dependent and reversible on one string, it isn't tied to a lexical class, and on Warren's account the pronoun may go either way. The QN override is none of these: near-categorical for core members, tied to a small conventional class, and construction-linked, as Appendix~\ref{sec:A-5}'s \mention{a number of} versus \mention{the number of} control shows on a single lexeme. Same surface phenomenon, different mechanism, and the frame's cells say which.

\subsection{Number-transparent NPs as a projectibility profile}\label{sec:2-4}

Section 1 set the secondary question: does identifying a positive access relation let us predict behaviour? For number-transparent quantificational nouns the profile hypothesis can be tested in two task settings: description of known members and class discovery for candidates. (The framework returns negatives too: the inanimate half of the \mention{bunch} contrast fails to project, §\ref{sec:2-2}, and coordination never qualifies, §\ref{sec:2-3}.)

For known \textit{CGEL} members (\mention{lot, plenty, lots, bags, heaps, loads, oodles, stacks, remainder, rest, number, couple}₂), lexical identity is sufficient input for description: knowing a noun is on this list licenses comparison with the rest of the profile. The non-trivial projective test is candidate discovery, where the input criteria cannot simply repeat the targets.

For candidate new members, class discovery uses two diagnostics independently observable from any of the projected behaviours: quantificational semantics (an indefinite quantity rather than a precise count or kind) and high-frequency occurrence as head of an \mention{of}-complement construction. A candidate that satisfies both is then tested against the projected cluster:

\begin{itemize}
\item
  near-obligatory agreement override in the central cases, singular- or plural-direction depending on subgroup;\footnote{A pilot COCA check found every fully audited discriminating cell in the predicted direction. Counter-direction cells received more extensive manual filtering than predicted-direction cells, so the counts are attestational corroboration rather than estimates of grammatical rates. \mention{A lot of money} required parse-shift filtering, and the large \mention{a lot of people is/was} cell remains incompletely scrubbed. Full counts and methodology: Appendix~\ref{sec:A-3}.}
\item
  characteristic determiner patterns by subgroup (\mention{lot} most often with \mention{a}, also with \mention{the} in totality uses; \mention{plenty} typically without a determiner; the plural-only forms typically without a determiner; \mention{remainder, rest} typically with \mention{the}; \mention{number, couple}₂ typically with \mention{a});
\item
  partitive vs.~pseudo-partitive \mention{of}-complement availability;
\item
  behaviour under complement ellipsis: the \mention{of}-complement can be omitted while remaining understood, and it continues to determine the NP's number (\mention{A lot was spent on travel} vs \mention{A lot were arrested}; \textit{CGEL} [56], p.~349);\footnote{These bare uses are not \textit{CGEL}'s \term{fused-head} construction, where the head fuses with a determiner or modifier (\mention{Did you buy \textup{[}some\textup{]} yesterday?}; §9.1, p.~410). \textit{CGEL} analyzes bare QN NPs as ordinary heads with an elliptical complement (pp.~412, 503). The fusion concept returns, differently applied, in the fused relatives of §\ref{sec:2-5}.}
\item
  count/non-count percolation of the \mention{of}-complement to the whole NP (\mention{lot} \enquote{allows the number of the oblique to percolate up to determine the number of the whole NP}; \textit{CGEL}, p.~349), with consequences for selection by predicates that are count- or non-count-sensitive;
\item
  restricted premodification (\mention{a whole/huge lot} but not \mention{a remarkable plenty}).
\end{itemize}

The class's sub-group structure makes the cluster's predictions concrete. Table~\ref{tab:qn-matrix} displays the per-noun pattern for the central diagnostics. The entries are \textit{CGEL}-coded except in the two \mention{of}-complement columns, where the COCA pilot has revised them: the determined frame is attested for every plural-only member, and \mention{rest} and \mention{remainder} admit a bare complement when it denotes uniquely. The semantic classification of \mention{rest} remains open, and a fuller per-noun corpus audit is the within-domain check still needed for the profile claim.

{
\begin{longtable}[]{@{}
  >{\raggedright\arraybackslash}p{(\linewidth - 8\tabcolsep) * \real{0.20}}
  >{\raggedright\arraybackslash}p{(\linewidth - 8\tabcolsep) * \real{0.16}}
  >{\raggedright\arraybackslash}p{(\linewidth - 8\tabcolsep) * \real{0.19}}
  >{\raggedright\arraybackslash}p{(\linewidth - 8\tabcolsep) * \real{0.20}}
  >{\raggedright\arraybackslash}p{(\linewidth - 8\tabcolsep) * \real{0.25}}@{}}
\caption{Number-transparent quantificational nouns: behavioural matrix (\textit{CGEL}-based pilot).}\label{tab:qn-matrix}\\
\toprule\noalign{}
Sub-group / Noun & Determiner & Partitive \mention{of} N & Pseudo-partitive \mention{of} N & Override direction \\
\midrule\noalign{}
\endfirsthead
\toprule\noalign{}
Sub-group / Noun & Determiner & Partitive \mention{of} N & Pseudo-partitive \mention{of} N & Override direction \\
\midrule\noalign{}
\endhead
\bottomrule\noalign{}
\endlastfoot
\multicolumn{5}{@{}l}{\textbf{Singular, \mention{a}-determined}} \\
\mention{lot} & \mention{a}; also \mention{the} (totality) & $\checkmark$ (\mention{a lot of the money}) & $\checkmark$ (\mention{a lot of money}) & sg → pl with plural complement \\
\mention{number} & \mention{a} & $\checkmark$ (plural oblique) & $\checkmark$ (plural oblique) & sg → pl with plural complement \\
\mention{couple}₂ & \mention{a} & $\checkmark$ (plural oblique) & $\checkmark$ (plural oblique) & sg → pl with plural complement \\
\addlinespace
\multicolumn{5}{@{}l}{\textbf{Singular, usually without a determiner}} \\
\mention{plenty} & usually without a determiner & not typical & $\checkmark$ (\mention{plenty of money}) & sg → pl with plural complement \\
\addlinespace
\multicolumn{5}{@{}l}{\textbf{Plural, usually without a determiner}} \\
\mention{lots, bags, heaps, loads, oodles, stacks} & usually without a determiner & $\checkmark$, dispreferred; attested for every member (§\ref{sec:2-1}) & $\checkmark$ & pl → sg with non-count complement \\
\addlinespace
\multicolumn{5}{@{}l}{\textbf{Singular, \mention{the}-determined}} \\
\mention{remainder} & \mention{the} & $\checkmark$ (\mention{the remainder of the time}); bare whole only if bounded and unique (\ungram{\mention{the remainder of time}}) & n/a: not a quantity noun & sg/pl (complement-controlled) \\
\mention{rest} & \mention{the} & $\checkmark$ (\mention{the rest of the time}); bare whole only if bounded and unique (\mention{the rest of society}) & ? (classification disputed) & sg/pl (complement-controlled) \\
\addlinespace
\multicolumn{5}{@{}l}{\textbf{Boundary}} \\
\mention{majority, minority} & \mention{a} & $\checkmark$ & ? & variable \\
\mention{bunch} & \mention{a} & $\checkmark$ & $\checkmark$ & variable \\
\end{longtable}
}

One sentence per central sub-group shows the override directions the table codes.

\ea\label{ex:subgroup-sentences}
\ea \mention{A number of spots have appeared.}
\ex \mention{Plenty of people were stranded.}
\ex \mention{Heaps of money has been spent.}
\ex \mention{The rest of the money was returned.}
\z
\z

(\ref{ex:subgroup-sentences}a) and (\ref{ex:subgroup-sentences}c) repeat (\ref{ex:cgel-two-direction}): a singular \mention{a}-determined noun driven plural by its complement, a plural-only noun driven singular. (\ref{ex:subgroup-sentences}b) shows determinerless \mention{plenty} with the plural override; (\ref{ex:subgroup-sentences}d) shows the \mention{the}-determined sub-group, where direction is complement-controlled (plural instead with \mention{the rest of the people}). The directions in all four match the COCA pilot's counts in Table~\ref{tab:override-audit}.

\textit{CGEL} treats override as obligatory in the central cases; the pilot corroborates the direction in selected cells but does not audit every class member or estimate rates. Complement-ellipsis behaviour and count/non-count percolation remain \textit{CGEL}-based predictions rather than per-noun corpus results. Premodification is restricted across the class, though the restriction looks like one on modifier type rather than on premodification as such: degree and size modifiers are available (\mention{a whole/huge lot}; \mention{huge}, \mention{enormous}, and \mention{vast piles of money}), while evaluative ones are marginal (\mention{a remarkable plenty}), and \mention{plenty} takes minimal modification generally.

Each sub-group covaries: knowing a candidate noun's determiner pattern and partitive availability narrows the sub-group, which then licenses the override-direction prediction (and vice versa). One qualification on the partitive diagnostic. For \mention{remainder}, the absence of a quantity reading follows from what kind of noun it is; for \mention{rest}, the lexical-semantic classification remains unresolved. The diagnostic therefore carries less weight than the determiner and override behaviours (Appendix~\ref{a.2-partitive-vs.-pseudo-partitive-of-complements}). The boundary cases have weaker covariance, which is why §\ref{sec:2-2}'s gradient framing applies to them.

Table~\ref{tab:override-audit} displays the raw agreement-direction counts. The discriminating cells are the informative ones, where the head's number and the complement's diverge. The matched cells confirm the expected direction but cannot distinguish the override from head-driven agreement. Counter-direction cells received more extensive manual checking than predicted-direction cells, so the table reports an audit trail rather than estimated rates. The large \mention{a lot of people} cell remains asymmetric, and the smallest cells remain feasibility checks. Referential \mention{the number of} is a negative control: the same noun agrees head-driven under \mention{the} and overrides only under quantificational \mention{a} (Appendix~\ref{sec:A-5}), so the contrast belongs to the construction, not the lexeme.

\begin{table}
\centering
\small
\caption{Raw agreement-direction counts in the COCA pilot. Counts are after the filtering described in Appendix~\ref{sec:A-3}; the audit-status column records the evidential limitation.}\label{tab:override-audit}
\begin{tabular}{@{}p{0.22\linewidth}p{0.12\linewidth}rrp{0.38\linewidth}@{}}
\toprule
Subject & Predicted & Plural & Singular & Audit status \\
\midrule
\mention{a lot of people} & plural & 4,195 & 85 & Large counter-direction cell not symmetrically scrubbed \\
\mention{a lot of money} & singular & 1 & 90 & Counter-direction cell checked; parse-shift filtering \\
\mention{a number of people} & plural & 100 & 0 & False positives checked \\
\mention{lots of people} & plural & 348 & 0 & False positives/non-standard uses checked \\
\mention{lots of money} & singular & 0 & 24 & Small cell; false positive checked \\
\mention{plenty of people} & plural & 79 & 0 & Small cell \\
\mention{plenty of money} & singular & 0 & 6 & Small cell \\
\mention{the rest of the people} & plural & 19 & 0 & Small cell \\
\mention{the rest of the money} & singular & 0 & 19 & Small cell \\
\bottomrule
\end{tabular}
\end{table}

This is the prospective design for a new-member case: quantificational semantics and \mention{of}-complement-head distribution are inputs, while the cluster of agreement, determiner, partitive, ellipsis, count/non-count, and premodification behaviours supplies the targets. Boundary cases like \mention{majority} and \mention{minority} (§\ref{sec:2-2}) and the dual-sense \mention{couple}₁ (collective) versus \mention{couple}₂ (number-transparent) show graded membership. The design distinguishes prediction from redescription, but the present evidence supports only selected targets.

The more demanding test asks what the profile predicts beyond the in-grammar cluster. Diachronic entry and cross-linguistic analogues are plausible targets, but neither is tested here. One candidate-member test has been run. \mention{Piles} was admitted on two criteria only, quantity meaning and frequent \mention{of}-complement distribution, neither of which is among the predicted behaviours. The determiner-frame prediction receives partial support: \mention{a piles of} returns no COCA hits, as the plural-only subgroup requires, while the determined frame \mention{the piles of} N selects the literal reading in the available sample. This result separates the two senses of \mention{piles}, but it does not validate the full profile.

A follow-up sweep suggests that the separation extends across the plural-only subgroup. The determined frame recruits the literal homonym wherever a member has one: \mention{stacks of the *} is dominated by library stacks and smokestacks, \mention{heaps of the *} by literal and colliery heaps, and \mention{bags} and \mention{loads} by containers and cargo. \mention{Oodles}, the one member with no literal homonym, returns quantificational hits exclusively. This is a useful subgroup hypothesis, not a completed projection test.

The agreement prediction couldn't be tested at all. COCA has 91 tokens of \mention{piles of money} and one followed by a finite form of \mention{be}; adding COHA's two centuries brings the discriminating non-count cell to three tokens, one singular and two plural. Nothing there supports the override and three tokens cannot weigh against it either. A container-noun control (\mention{sacks/boxes/crates of money}) returns nothing at all, and the left contexts of the 91 tokens locate the problem precisely: 41 have \mention{piles of money} as the object of a verb (\mention{make, spend, earn, amass, inherit}) and 13 as the complement of a preposition, so at least 59\% sit in complement position, while the single token preceded by finite \mention{be} is copular or existential, with agreement running from a different subject. Quantity expressions of this shape occur as complements, not as subjects of finite \mention{be}. The same sparsity blocked the inanimate \mention{bunch} contrast in §\ref{sec:2-2}.

So the override prediction, which is the profile's headline, is corpus-testable only for members that appear in subject position, and candidate-member discovery for this class will need acceptability judgments rather than frequency. The \mention{piles} result is therefore partial prospective support plus a feasibility failure. The protocol, full returns, and corrections the wildcard queries forced on Appendix~\ref{a.2-partitive-vs.-pseudo-partitive-of-complements} are in the supplementary material (see Data availability).

The cluster makes number-transparent quantificational nouns the strongest candidate \term{projectibility profile} examined here: an observable base supports projection to a structured set of targets, with conventional lexical inheritance as the proposed stabilizer. The profile is a candidate for Slater's \autocite*{slater-2015-natural-kindness} stability-plus-projection diagnostic, but the within-domain check beyond selected agreement and determiner behaviour remains open. A profile prediction would fail if a noun satisfying the input diagnostics systematically lacked the target cluster. The status of \mention{rest} is an explicit analytical fork: it remains an access-transparency case under either treatment, but if it is a part/subset expression rather than a quantity noun, it should not count as input to the quantity-QN candidate-discovery generalization. Whether the profile earns a local-kind label is a downstream question for §\ref{sec:7}.

\subsection{Transparent free relatives}\label{sec:2-5}

A second morphosyntactic-transparency case fits the schema. Transparent free relatives (TFRs) are a sub-type of fused-relative construction (\textit{CGEL}, ch.~12, §6, pp.~1068ff.) in which the matrix accesses properties of an embedded predicate (the \term{nucleus}) through \mention{what}: \mention{what they call ecstatic}, \mention{what we regard as essential}, \mention{what seems to be a tourist}. \term{Fused} is \textit{CGEL}'s term for one word jointly realizing two functions: in a fused relative, \enquote{the head of an NP fuses with the relativised element in a relative clause} (p.~1073), so \mention{what} in \mention{what she wrote} is at once the NP's head and the relativized element. This paper uses TFRs only as a second profile case. The analysis followed here is Kim's construction-grammar one, developed across \textcite{kim-2011-tfrs}, \textcite{kim-2017-free-relatives}, \textcite[ch.~6.8]{kim-2026-form-function}, and its revision in \textcite[§5]{kim-2026-indeterminacy-free-relatives}, with the corpus material in \textcite{kim-reynolds-2026-tfr-two-kinds}. The verdict below is conditional on that analysis: rival treatments locate the mediator elsewhere, and where the mediator sits is exactly what the schema needs fixed, so a reader who prefers Wilder's \autocite*{wilder-1999-transparent-free-relatives} parenthetical-with-backward-deletion account, van Riemsdijk's \autocite*{vanriemsdijk-2006-free-relatives} graft analysis, or Grosu's \autocite*{grosu-2003-unified-theory-tfrs} unified treatment should expect the slots to be filled differently: each locates the mediator, and so the access path, somewhere else. The verdict is therefore analysis-relative: positive under Kim's constructional analysis, indeterminate on an analysis-neutral reading, since the competing treatments aren't adjudicated here.

\ea\label{ex:tfr-r-relations}
\ea \mention{What appear to be trousers are on the chair.}
\ex \mention{What they call essential workers are exhausted.}
\ex \mention{He was utterly penniless and what they call dimwitted.}
\z
\z

TFRs provide a second, conditional grammatical illustration, but the evidence is deliberately bounded. They support two relations at different levels of the construction. One is agreement: matrix agreement accesses the nucleus's number profile through singular \mention{what}. The other is matrix selection, whose accessible property is the nucleus's syntactic category; \term{category flexibility} is the outcome, that a non-NP nucleus can satisfy the matrix's requirement, as in (\ref{ex:tfr-r-relations}c), where the nucleus is an AP. Predicative position admits both NPs and APs, so that example shows the nucleus's category rather than establishing a category-discriminating environment. Kim \autocite*[36--37]{kim-2026-indeterminacy-free-relatives} treats the second as the marked sub-case of a wider \term{selection transparency}, on which the matrix's selectional requirement is satisfied by the nucleus generally, NP nuclei included, as when a preposition selects the nucleus through the \mention{what}-clause (\mention{he speaks in what linguists call a Northern dialect}).

The concealing counterpart for parent-level agreement is the ordinary fused relative, where \mention{what} itself controls agreement. That control is family-level rather than minimally matched: an ordinary fused relative like \mention{what she wrote} lacks a nucleus, so it can't hold the concealed property constant while varying only the access. It establishes that the fused-relative family admits mediator-controlled agreement, not that a matched variant conceals this nucleus's number. The fourth gate requires the counterpart's evidential strength to be stated, and it enters here at that weaker, family-level grade.

The pilot reported with the joint analysis supplies a bounded check on the category contrast. In a full coding of 485 archived \mention{call} hits, 185 were genuine AP-nucleus TFRs, 144 of them externally predicative, with one token coordinated into attributive position. Exhaustive checks of the candidate sets for right-edge non-NP nuclei with \mention{seem} or \mention{appear} (721 tokens across eight queries, run or re-verified against COCA in August 2026) found no adjectival or adverbial case and a single PP-nucleus token in supplement position; the working paper reports the query details. This is distributional evidence, not a completed acceptability study; the relevant category-discriminating environments and judgments remain to be added to the joint paper.

The observable base is constructional sub-type plus verb-class membership, and the examples instantiate both classes Kim \autocite*[40]{kim-2026-indeterminacy-free-relatives} distinguishes: the \term{evidential} type (\mention{seem~to~be, appear~to~be, happen~to~be}; \ref{ex:tfr-r-relations}a), which takes NP nuclei only in the current analysis, and the \term{attributional} type (\mention{call, consider, describe~as, regard~as, take~to~be}; \ref{ex:tfr-r-relations}b,~c), which licenses AP, AdvP, PP, and \mention{-ing} nuclei as well. On the inheritance hierarchy Kim \autocite*[46]{kim-2026-indeterminacy-free-relatives} sets out, the parent licenses agreement, selection, binding, and coordination transparency across both classes, and the attributional daughter adds category flexibility. \enquote{Coordination transparency} there is Kim's name for a TFR diagnostic, that the construction conjoins as its nucleus's category; it doesn't reverse §\ref{sec:2-3}'s verdict on ordinary coordinate NPs, which concerns a different configuration. The pilot therefore supports a provisional parent/daughter contrast, with constructional inheritance as the proposed stabilizer. The full verb-class inventory, AP/PP stratification, and rival analyses remain part of the joint paper.

\subsection{Bridge}

The cases in Section \ref{sec:2} share a relation: morphosyntactic matching mediated by a construction or lexical class. Section \ref{sec:3} marks form--meaning applications at the boundary; the cross-linguistic reach of the schema, including the controlled Spanish and Tsez cases and the topology survey, is Section \ref{sec:6}'s topic.
